\documentclass[11pt,a4paper]{article}

\usepackage[indonesian]{babel}
\usepackage[T1]{fontenc}
\usepackage{lmodern}
\usepackage{amsmath}
\usepackage[margin=2.5cm]{geometry}
\usepackage{booktabs,tabularx,longtable,array,multirow}
\usepackage{xcolor,graphicx,float,caption}
\usepackage{enumitem}
\usepackage{listings}
\usepackage{algorithm}
\usepackage{algpseudocode}
\usepackage{tikz}
\usetikzlibrary{positioning,arrows.meta}
\usepackage[hidelinks]{hyperref}

\setlength{\parskip}{0.5em}
\setlength{\parindent}{0pt}
\renewcommand{\arraystretch}{1.25}
\newcolumntype{L}{>{\raggedright\arraybackslash}X}
\newcolumntype{P}[1]{>{\raggedright\arraybackslash}p{#1}}
\newcommand{\cond}[1]{\texttt{#1}}

\floatname{algorithm}{Algoritma}
\renewcommand{\algorithmicrequire}{\textbf{Masukan:}}
\renewcommand{\algorithmicensure}{\textbf{Keluaran:}}
\algrenewcommand\algorithmicif{\textbf{if}}
\algrenewcommand\algorithmicthen{\textbf{then}}

\lstset{
  basicstyle=\ttfamily\footnotesize,
  frame=single,
  breaklines=true,
  columns=fullflexible,
  keepspaces=true,
}

% ------------------------------------------------------------------
% TODO: lengkapi nama anggota dan mata kuliah.
\title{\textbf{Sistem Pakar Rekomendasi Berbasis\\Ripple Down Rules (RDR)}\\[0.5em]
\large Laporan Tugas \textless Nama Mata Kuliah\textgreater}
\author{
  \textless Nama Anggota 1\textgreater{} (13523065)\\
  \textless Nama Anggota 2\textgreater{} (13523117)
}
\date{\today}
% ------------------------------------------------------------------

\begin{document}
\maketitle
\tableofcontents
\newpage

% ==================================================================
\section{Pendahuluan}

\subsection{Deskripsi Tugas}
Tugas ini meminta pembuatan sistem AI berbasis \emph{Ripple Down Rules} (RDR) yang dipakai
seorang pakar untuk memberi rekomendasi kepada konsumen. Ketentuan utamanya:
\begin{enumerate}[nosep]
  \item Knowledge base (KB) awal hanya berisi satu rule, $t \rightarrow$ kesimpulan default,
        yang dapat diisi sesuai domain masalah.
  \item Kesimpulan diperoleh dengan menelusuri binary tree. Anak kanan (\emph{true-child})
        adalah pengecualian yang dieksekusi jika kondisi node terpenuhi; anak kiri
        (\emph{false-child}) dieksekusi jika kondisi tidak terpenuhi.
  \item Kesimpulan yang diberikan adalah kesimpulan terakhir yang diperoleh selama penelusuran.
  \item Jika pakar setuju, kesimpulan diberikan dan proses selesai. Jika tidak, dibuat rule
        pengecualian baru dengan kesimpulan dari pakar. Kondisinya dipilih dari kondisi yang
        dipenuhi kasus baru tetapi tidak dipenuhi kasus kunci (\emph{cornerstone case}) dari
        kesimpulan yang dianulir.
  \item Atribut dapat berjenis numerik maupun kategorikal.
\end{enumerate}

\subsection{Ringkasan Solusi}
Kami membangun engine \emph{Single Classification RDR} (SCRDR) yang generik terhadap domain.
Domain (nama, kesimpulan default, dan daftar atribut) dibaca dari file konfigurasi, sehingga
sistem yang sama dapat dipakai untuk domain apa pun. Engine dilengkapi dua antarmuka:
aplikasi web lokal (Streamlit) dengan visualisasi tree, dan mode terminal (CLI) yang dapat
merekam sesi ke file teks. Domain contoh yang dipakai pada laporan ini adalah
\emph{rekomendasi produk keuangan}.

% ==================================================================
\section{Dasar Teori}

\subsection{Ripple Down Rules}
Ripple Down Rules adalah metode akuisisi pengetahuan yang diperkenalkan oleh Compton dan
Jansen. Alih-alih meminta pakar menyusun seluruh basis pengetahuan di awal, pengetahuan
ditambahkan secara bertahap ketika sistem salah menangani suatu kasus. Pakar cukup
menjelaskan \emph{apa yang membedakan} kasus baru dari kasus lama yang menjadi alasan rule
sebelumnya dibuat. Pengetahuan lama tidak pernah diubah atau dihapus, hanya diberi
pengecualian.

\subsection{Struktur Single Classification RDR}
Pada SCRDR, KB berbentuk binary tree. Setiap node menyimpan:
\begin{itemize}[nosep]
  \item \textbf{kondisi}: konjungsi kondisi atomik (root memiliki kondisi \emph{true});
  \item \textbf{kesimpulan};
  \item \textbf{cornerstone case}: kasus yang menyebabkan rule tersebut dibuat;
  \item \textbf{true-child} (kanan): pengecualian, dikunjungi jika kondisi TRUE;
  \item \textbf{false-child} (kiri): alternatif, dikunjungi jika kondisi FALSE.
\end{itemize}

Karena kondisi root selalu TRUE, false-child root tidak pernah dikunjungi. Seluruh rule yang
menganulir kesimpulan default membentuk rantai false-child di bawah true-child root, yang dapat
dibaca sebagai ``default, \emph{kecuali} $a \rightarrow X$, \emph{atau jika tidak}
$e \rightarrow Y$, \ldots''.

\subsection{Inferensi}
Penelusuran dimulai dari root. Jika kondisi node terpenuhi, kesimpulannya dicatat lalu
penelusuran berlanjut ke true-child; jika tidak, ke false-child. Penelusuran berhenti saat
child yang dituju kosong. Kesimpulan akhir adalah kesimpulan node terakhir yang bernilai TRUE.

\subsection{Akuisisi Pengetahuan dan Difference List}
Saat pakar tidak setuju, dua node dari hasil penelusuran menjadi penting:
\begin{itemize}[nosep]
  \item \textbf{lastTrue}: node terakhir bernilai TRUE, yaitu pemberi kesimpulan yang salah.
        \emph{Cornerstone} node ini menjadi pembanding.
  \item \textbf{lastEval}: node terakhir yang dievaluasi. Node baru ditempel di sini, sebagai
        true-child jika kondisinya TRUE atau false-child jika FALSE.
\end{itemize}
\emph{Difference list} adalah himpunan kondisi yang bernilai TRUE untuk kasus baru dan FALSE
untuk cornerstone \textbf{lastTrue}. Dengan memilih kondisi dari daftar ini, cornerstone lama
dijamin tetap mendapat kesimpulan lamanya, sehingga tidak terjadi regresi.

Cornerstone \textbf{lastTrue} selalu mencapai posisi node baru: semua pengecualian
\textbf{lastTrue} yang sudah ada dibuat agar FALSE untuk cornerstone tersebut, sehingga
cornerstone itu menelusuri rantai pengecualian sampai ujung. Karena itu cukup cornerstone ini
saja yang perlu dibedakan.

% ==================================================================
\section{Analisis dan Desain}

\subsection{Arsitektur}
Sistem dibagi menjadi tiga lapis yang terpisah:
\begin{enumerate}[nosep]
  \item \textbf{Engine} (\texttt{rdr/}): struktur data, inferensi, difference list, penambahan
        rule, validasi, serta penyimpanan. Engine tidak bergantung pada antarmuka apa pun.
  \item \textbf{Antarmuka}: aplikasi web Streamlit (\texttt{app.py}, \texttt{ui/}) dan mode
        terminal (\texttt{cli.py}).
  \item \textbf{Data}: satu folder per domain di \texttt{domains/}, berisi
        \texttt{domain.json} (ditulis manusia) dan \texttt{kb.json} (ditulis sistem).
\end{enumerate}

\subsection{Struktur Data}
\begin{table}[H]
\centering
\begin{tabularx}{\textwidth}{lL}
\toprule
\textbf{Struktur} & \textbf{Isi} \\
\midrule
\texttt{Case} & Pemetaan nama atribut ke nilai (float untuk numerik, string untuk kategorikal). \\
\texttt{Condition} & Tripel (atribut, operator, nilai). Operator numerik: \cond{>}, \cond{>=},
  \cond{<}, \cond{<=}. Operator kategorikal: \cond{==}, \cond{!=}. Atribut yang tidak ada pada
  kasus membuat kondisi bernilai FALSE. \\
\texttt{Node} & id, daftar kondisi (AND), kesimpulan, cornerstone, true-child, false-child, serta
  metadata posisi (parent, cabang) dan id node yang dianulir. \\
\texttt{KnowledgeBase} & Root dan penghitung id; operasi \texttt{infer}, \texttt{add\_rule},
  \texttt{difference\_list}, \texttt{verify}. \\
\texttt{Domain} & Nama, kesimpulan default, dan daftar \texttt{Attribute} (nama, tipe, daftar
  nilai opsional, min/max/satuan opsional). \\
\bottomrule
\end{tabularx}
\caption{Struktur data utama}
\end{table}

\subsection{Konfigurasi Domain}
Domain didefinisikan oleh \texttt{domain.json}. Informasi minimal yang diperlukan adalah
kesimpulan default serta nama dan tipe setiap atribut. Tipe atribut diperlukan untuk
menentukan operator kandidat dan cara membandingkan nilai.

\begin{lstlisting}
{
  "name": "Rekomendasi Produk Keuangan",
  "default_conclusion": "Tabungan Reguler",
  "attributes": {
    "umur":   {"type": "numeric", "min": 17, "max": 100, "unit": "tahun"},
    "gaji":   {"type": "numeric", "min": 0, "unit": "juta/bulan"},
    "risiko": {"type": "categorical", "values": ["rendah", "sedang", "tinggi"]},
    "status": {"type": "categorical", "values": ["pelajar", "pekerja", "pensiun"]}
  }
}
\end{lstlisting}

Domain baru dapat dibuat lewat antarmuka (formulir di Streamlit atau wizard di CLI) maupun
dengan menulis \texttt{domain.json} secara manual. Atribut baru dapat ditambahkan kapan saja.

\subsection{Keputusan Desain dan Asumsi}
Beberapa hal tidak dirinci oleh spesifikasi. Tabel~\ref{tab:asumsi} merangkum keputusan kami.

\begin{longtable}{P{0.3\textwidth}P{0.64\textwidth}}
\toprule
\textbf{Aspek} & \textbf{Keputusan} \\
\midrule
\endhead
Jumlah kondisi per rule & Boleh lebih dari satu (digabung AND). Setiap kondisi harus berasal
  dari difference list. \\
Pemilih kondisi & Pakar memilih dari difference list yang ditampilkan sistem. Sistem
  memvalidasi pilihan. \\
Kandidat kondisi numerik & Relatif terhadap cornerstone. Jika nilai kasus baru $v > c$
  (nilai cornerstone): \cond{attr > c} (paling umum) dan \cond{attr >= v} (paling spesifik).
  Jika $v < c$: \cond{attr < c} dan \cond{attr <= v}. Jika $v = c$: tidak ada kandidat.
  Tanpa cornerstone: \cond{attr >= v} dan \cond{attr <= v}. \\
Kandidat kondisi kategorikal & \cond{attr == v}, ditambah \cond{attr != c} jika nilai
  cornerstone berbeda. \\
Threshold kustom & Pakar boleh memasukkan kondisi sendiri (misalnya \cond{umur > 45}) selama
  kondisi itu TRUE untuk kasus baru dan FALSE untuk cornerstone. \\
Cornerstone root & Root tidak memiliki cornerstone, sehingga rule yang menganulir default boleh
  memakai kondisi apa pun dari kasus baru. \\
Difference list kosong & Disebut \emph{conflict}. Rule lama tidak diubah. Pakar dapat membatalkan
  koreksi atau menambah atribut baru yang membedakan kedua kasus. \\
Atribut baru & Cornerstone lama tidak memiliki nilai atribut tersebut, sehingga kondisi pada
  atribut itu bernilai FALSE untuk cornerstone lama. \\
Nilai kosong & Semua atribut wajib diisi saat konsultasi. \\
Kesimpulan sama & Jika pakar memberi kesimpulan yang sama dengan sistem, dianggap setuju dan
  tidak ada rule baru. \\
Penulisan kesimpulan & Dinormalisasi (spasi, huruf besar/kecil) mengikuti kesimpulan yang
  sudah ada, misalnya ``saham'' menjadi ``Saham''. \\
Persistensi & KB disimpan otomatis ke \texttt{kb.json} setiap kali rule ditambahkan. Reset
  membuat backup file lama. \\
\bottomrule
\caption{Keputusan desain dan asumsi}
\label{tab:asumsi}
\end{longtable}

% ==================================================================
\section{Algoritma}

\begin{algorithm}[H]
\caption{Inferensi}
\begin{algorithmic}[1]
\Require root KB, kasus $x$
\Ensure kesimpulan, lastTrue, lastEval, lastResult, path
\State $node \gets root$; \quad $lastTrue \gets root$; \quad $lastEval \gets root$
\State $lastResult \gets$ \textbf{true}; \quad $path \gets [\,]$
\While{$node \neq \textbf{null}$}
  \State $lastEval \gets node$
  \If{semua kondisi $node$ dipenuhi $x$}
    \State $lastTrue \gets node$; \quad $lastResult \gets$ \textbf{true}
    \State tambahkan $(node, \textbf{true})$ ke $path$
    \State $node \gets node.trueChild$ \Comment{anak kanan: pengecualian}
  \Else
    \State $lastResult \gets$ \textbf{false}
    \State tambahkan $(node, \textbf{false})$ ke $path$
    \State $node \gets node.falseChild$ \Comment{anak kiri: alternatif}
  \EndIf
\EndWhile
\State \Return $(lastTrue.conclusion, lastTrue, lastEval, lastResult, path)$
\end{algorithmic}
\end{algorithm}

\begin{algorithm}[H]
\caption{Pembangkitan Difference List}
\begin{algorithmic}[1]
\Require kasus baru $x$, cornerstone $c$ (boleh \textbf{null}), domain
\Ensure daftar kondisi pembeda
\State $cand \gets [\,]$
\For{setiap atribut $a$ pada domain}
  \State $v \gets x[a]$; \quad $w \gets c[a]$ jika ada, selain itu \textbf{null}
  \If{$a$ numerik}
    \If{$w = \textbf{null}$} tambahkan $a \geq v$, $a \leq v$
    \ElsIf{$v > w$} tambahkan $a > w$, $a \geq v$
    \ElsIf{$v < w$} tambahkan $a < w$, $a \leq v$
    \EndIf
  \Else
    \State tambahkan $a = v$
    \If{$w \neq \textbf{null}$ \textbf{and} $w \neq v$} tambahkan $a \neq w$ \EndIf
  \EndIf
\EndFor
\State \Return $[\,k \in cand \mid k(x) = \textbf{true} \wedge k(c) = \textbf{false}\,]$
\end{algorithmic}
\end{algorithm}

\begin{algorithm}[H]
\caption{Penambahan Rule}
\begin{algorithmic}[1]
\Require KB, kasus $x$, kesimpulan pakar $k$, kondisi terpilih $C$
\State $(concl, lastTrue, lastEval, lastResult, \_) \gets \Call{Infer}{root, x}$
\If{$k = concl$} \textbf{tolak}: kesimpulan sama \EndIf
\If{$C$ kosong} \textbf{tolak}: minimal satu kondisi \EndIf
\For{setiap $q \in C$}
  \If{$q(x) = \textbf{false}$} \textbf{tolak}: tidak dipenuhi kasus baru \EndIf
  \If{$q(lastTrue.cornerstone) = \textbf{true}$} \textbf{tolak}: cornerstone ikut teranulir \EndIf
\EndFor
\State $n \gets \Call{Node}{C, k, cornerstone = \text{salinan}(x)}$
\If{$lastResult$}
  \State $lastEval.trueChild \gets n$
\Else
  \State $lastEval.falseChild \gets n$
\EndIf
\State simpan KB
\end{algorithmic}
\end{algorithm}

\begin{algorithm}[H]
\caption{Siklus Konsultasi}
\begin{algorithmic}[1]
\State muat domain dan KB (jika belum ada: $root \gets$ IF true THEN default)
\Loop
  \State $x \gets$ input data konsumen (divalidasi sesuai tipe atribut)
  \State $res \gets \Call{Infer}{root, x}$; tampilkan kesimpulan dan jalur
  \If{pakar setuju}
    \State berikan $res.conclusion$ kepada konsumen
  \Else
    \State $k \gets$ kesimpulan yang benar dari pakar
    \State $D \gets \Call{DifferenceList}{x, res.lastTrue.cornerstone}$
    \While{$D$ kosong} \Comment{conflict}
      \State pakar membatalkan, \textbf{atau} menambah atribut baru dan nilainya untuk $x$
      \State $D \gets \Call{DifferenceList}{x, res.lastTrue.cornerstone}$
    \EndWhile
    \State $C \gets$ kondisi yang dipilih pakar dari $D$
    \State \Call{AddRule}{x, k, C}; berikan $k$ kepada konsumen
  \EndIf
\EndLoop
\end{algorithmic}
\end{algorithm}

% ==================================================================
\section{Implementasi}

\subsection{Struktur Proyek}
\begin{lstlisting}
app.py               entry point Streamlit
cli.py               mode terminal
rdr/                 engine (tidak bergantung pada UI)
  models.py          Condition, Node
  engine.py          KnowledgeBase: infer(), add_rule(), verify()
  difference.py      kandidat kondisi dan difference list
  schema.py          Domain, Attribute, validasi kasus, load/save domain.json
  storage.py         load/save/reset kb.json
ui/                  halaman Streamlit
domains/             satu folder per domain (domain.json, kb.json)
scripts/             skenario demo
tests/               pytest
\end{lstlisting}

\subsection{Antarmuka Web (Streamlit)}
Aplikasi web terdiri dari empat tab:
\begin{itemize}[nosep]
  \item \textbf{Konsultasi}: formulir data konsumen yang dibangun otomatis dari schema domain.
        Menampilkan rekomendasi, tabel jalur penelusuran, dan posisi kasus pada tree. Dalam mode
        pakar tersedia tombol setuju/tidak setuju, pilihan kesimpulan, tabel perbandingan
        kasus baru dengan cornerstone, difference list, kondisi kustom, pratinjau rule beserta
        posisinya, dan penanganan conflict.
  \item \textbf{Tree \& Rule}: visualisasi binary tree (anak kiri F berwarna merah, anak kanan
        T berwarna hijau, jalur kasus terakhir disorot), daftar rule beserta cornerstone,
        verifikasi KB, unduh \texttt{kb.json}, dan reset KB.
  \item \textbf{Uji Batch (CSV)}: menjalankan inferensi untuk banyak kasus sekaligus.
  \item \textbf{Domain \& Atribut}: melihat schema, mengubah default, menambah atribut, dan
        membuat domain baru.
\end{itemize}
Sidebar berisi pemilihan domain dan toggle \emph{mode pakar}. Ketika toggle dimatikan, sistem
berperan sebagai mode konsumen yang hanya menampilkan rekomendasi.

% TODO: ganti kotak di bawah dengan screenshot aplikasi (simpan di images/).
\begin{figure}[H]
\centering
\fbox{\parbox[c][5cm][c]{0.85\textwidth}{\centering Screenshot tab Konsultasi\\
  (\texttt{images/konsultasi.png})}}
\caption{Tab Konsultasi dalam mode pakar}
\end{figure}

\begin{figure}[H]
\centering
\fbox{\parbox[c][5cm][c]{0.85\textwidth}{\centering Screenshot tab Tree \& Rule\\
  (\texttt{images/tree.png})}}
\caption{Visualisasi tree pada tab Tree \& Rule}
\end{figure}

\subsection{Mode Terminal (CLI)}
Mode terminal menyediakan alur yang sama melalui menu teks: konsultasi, koreksi pakar,
difference list bernomor, kondisi kustom, penanganan conflict, tree ASCII, daftar rule,
verifikasi, penambahan atribut, serta pembuatan domain. Opsi \texttt{--log} merekam seluruh
prompt dan jawaban ke file teks. Potongan rekaman sesi:

\begin{lstlisting}
Data konsumen:
  umur (angka, 17..100, tahun): 60
  gaji (angka, 0..inf, juta/bulan): 20
  risiko (rendah/sedang/tinggi): tinggi
  status (pelajar/pekerja/pensiun): pensiun

Jalur penelusuran:
  R0: IF true THEN Tabungan Reguler        -> TRUE (ke anak kanan)
  R1: IF risiko == tinggi THEN Saham       -> TRUE (ke anak kanan)

REKOMENDASI: Saham   (dari R1)
Pakar setuju? [y/n]: n
Kesimpulan yang benar (ketik nomor atau teks baru, kosong = batal):
  1. Tabungan Reguler
> Deposito

Kesimpulan 'Saham' (R1) akan dianulir.
  Atribut | Kasus baru | Cornerstone R1
  --------+------------+---------------
  umur    | 60         | 30
  gaji    | 20         | 15
  risiko  | tinggi     | tinggi
  status  | pensiun    | pekerja

Difference list:
  1. umur > 30
  2. umur >= 60
  3. gaji > 15
  4. gaji >= 20
  5. status == pensiun
  6. status != pekerja
> 5
  Kondisi: status == pensiun
  Sudah cukup? [y/n]: y

Rule R2 ditambahkan: IF status == pensiun THEN Deposito
Ditempel di R1.trueChild, pengecualian dari R1.
\end{lstlisting}

\subsection{Cara Menjalankan}
\begin{lstlisting}
pip install -r requirements.txt
streamlit run app.py                                   # aplikasi web
python cli.py --domain keuangan                        # mode terminal (pakar)
python cli.py --domain keuangan --consumer             # mode konsumen
python cli.py --domain keuangan --log docs/sesi.txt    # rekam sesi
python -m scripts.demo_scenario                        # jalankan skenario C1-C9
python -m pytest                                       # pengujian
\end{lstlisting}

% ==================================================================
\section{Pengujian}

\subsection{Domain Uji}
Domain rekomendasi produk keuangan memiliki kesimpulan default \emph{Tabungan Reguler} dan
empat atribut: \texttt{umur} (numerik, tahun), \texttt{gaji} (numerik, juta/bulan),
\texttt{risiko} (kategorikal: rendah, sedang, tinggi), dan \texttt{status} (kategorikal:
pelajar, pekerja, pensiun). Pengujian dimulai dari KB yang hanya berisi
$R0$: IF true THEN Tabungan Reguler.

\subsection{Skenario Input dan Output}
Tabel~\ref{tab:skenario} menunjukkan sembilan kasus yang dimasukkan secara berurutan. Notasi
$R_i$(T) berarti kondisi $R_i$ terpenuhi sehingga penelusuran berlanjut ke anak kanan;
$R_i$(F) berarti tidak terpenuhi sehingga berlanjut ke anak kiri.

{\small
\begin{longtable}{P{0.5cm}P{2.7cm}P{3.1cm}P{1.9cm}P{1.6cm}P{3.4cm}}
\toprule
\textbf{\#} & \textbf{Input (umur, gaji, risiko, status)} & \textbf{Jalur} &
\textbf{Output sistem} & \textbf{Pakar} & \textbf{Aksi} \\
\midrule
\endhead
C1 & 20, 2, rendah, pelajar & R0(T) & Tabungan Reguler & Setuju & -- \\
C2 & 30, 15, tinggi, pekerja & R0(T) & Tabungan Reguler & Saham &
  R1 di R0.trueChild \\
C3 & 25, 8, sedang, pekerja & R0(T), R1(F) & Tabungan Reguler & Reksadana &
  R2 di R1.falseChild \\
C4 & 60, 20, tinggi, pensiun & R0(T), R1(T) & Saham & Deposito & R3 di R1.trueChild \\
C5 & 35, 25, tinggi, pekerja & R0(T), R1(T), R3(F) & Saham & Setuju & -- \\
C6 & 22, 3, tinggi, pelajar & R0(T), R1(T), R3(F) & Saham & Reksadana &
  R4 di R3.falseChild \\
C7 & 65, 2, tinggi, pensiun & R0(T), R1(T), R3(T) & Deposito & Setuju &
  -- (R4 tidak dicek walau \cond{gaji <= 3} benar) \\
C8 & 50, 40, sedang, pekerja & R0(T), R1(F), R2(T) & Reksadana & Obligasi &
  R5 di R2.trueChild \\
C9 & 30, 15, tinggi, pekerja & R0(T), R1(T), R3(F), R4(F) & Saham & Obligasi &
  \textbf{Conflict}: difference list kosong \\
\bottomrule
\caption{Skenario pengujian C1--C9}
\label{tab:skenario}
\end{longtable}
}

\subsection{Pembentukan Rule}
Tabel~\ref{tab:diff} menunjukkan difference list yang dihasilkan sistem untuk setiap koreksi
dan kondisi yang dipilih pakar.

{\small
\begin{longtable}{P{0.8cm}P{1.4cm}P{3.2cm}P{4.6cm}P{3.2cm}}
\toprule
\textbf{Rule} & \textbf{Kasus} & \textbf{Pembanding} & \textbf{Difference list} &
\textbf{Rule terbentuk} \\
\midrule
\endhead
R1 & C2 & R0 (tanpa cornerstone) & \cond{umur >= 30}, \cond{umur <= 30}, \cond{gaji >= 15},
  \cond{gaji <= 15}, \cond{risiko == tinggi}, \cond{status == pekerja} &
  IF \cond{risiko == tinggi} THEN Saham \\
R2 & C3 & R0 (tanpa cornerstone) & \cond{umur >= 25}, \cond{umur <= 25}, \cond{gaji >= 8},
  \cond{gaji <= 8}, \cond{risiko == sedang}, \cond{status == pekerja} &
  IF \cond{risiko == sedang} THEN Reksadana \\
R3 & C4 & C2 (cornerstone R1) & \cond{umur > 30}, \cond{umur >= 60}, \cond{gaji > 15},
  \cond{gaji >= 20}, \cond{status == pensiun}, \cond{status != pekerja} &
  IF \cond{status == pensiun} THEN Deposito \\
R4 & C6 & C2 (cornerstone R1) & \cond{umur < 30}, \cond{umur <= 22}, \cond{gaji < 15},
  \cond{gaji <= 3}, \cond{status == pelajar}, \cond{status != pekerja} &
  IF \cond{gaji <= 3} THEN Reksadana \\
R5 & C8 & C3 (cornerstone R2) & \cond{umur > 25}, \cond{umur >= 50}, \cond{gaji > 8},
  \cond{gaji >= 40} & IF \cond{gaji >= 40} THEN Obligasi \\
-- & C9 & C2 (cornerstone R1) & \emph{kosong} & Conflict \\
\bottomrule
\caption{Difference list dan rule yang terbentuk}
\label{tab:diff}
\end{longtable}
}

Dua hal penting dari tabel di atas:
\begin{itemize}[nosep]
  \item \textbf{R2} ditempel di bawah R1 (karena R1 adalah lastEval), tetapi yang dianulir
        adalah kesimpulan R0. Karena itu R2 dibandingkan dengan cornerstone R0, yang kosong.
  \item \textbf{R4} ditempel di bawah R3 (lastEval), tetapi dibandingkan dengan cornerstone R1
        (C2), bukan cornerstone R3. Cornerstone R3 (C4) bernilai TRUE di R3 sehingga tidak
        pernah mencapai false-child R3.
\end{itemize}

\subsection{Tree Akhir}
\begin{figure}[H]
\centering
\begin{tikzpicture}[
  every node/.style={draw, rounded corners, align=center, font=\footnotesize,
                     text width=2.7cm, inner sep=3pt, fill=gray!8},
  tedge/.style={-{Stealth}, draw=green!50!black, thick},
  fedge/.style={-{Stealth}, draw=red!70!black, thick},
  lbl/.style={draw=none, fill=none, text width=, font=\scriptsize\bfseries}
]
\node[text width=3.3cm] (r0) at (0, 0) {R0\\\texttt{true}\\$\rightarrow$ Tabungan Reguler};
\node (r1) at (3.4, -1.9)  {R1\\\texttt{risiko == tinggi}\\$\rightarrow$ Saham};
\node (r2) at (0, -3.9)    {R2\\\texttt{risiko == sedang}\\$\rightarrow$ Reksadana};
\node (r3) at (6.8, -3.9)  {R3\\\texttt{status == pensiun}\\$\rightarrow$ Deposito};
\node (r5) at (1.7, -5.9)  {R5\\\texttt{gaji >= 40}\\$\rightarrow$ Obligasi};
\node (r4) at (5.1, -5.9)  {R4\\\texttt{gaji <= 3}\\$\rightarrow$ Reksadana};
\draw[tedge] (r0) -- node[lbl, above right, text=green!50!black] {T} (r1);
\draw[fedge] (r1) -- node[lbl, above left, text=red!70!black] {F} (r2);
\draw[tedge] (r1) -- node[lbl, above right, text=green!50!black] {T} (r3);
\draw[tedge] (r2) -- node[lbl, above right, text=green!50!black] {T} (r5);
\draw[fedge] (r3) -- node[lbl, above left, text=red!70!black] {F} (r4);
\end{tikzpicture}
\caption{Tree akhir setelah skenario C1--C9 (kiri = F/false-child, kanan = T/true-child)}
\end{figure}

\subsection{Uji Regresi Cornerstone}
Setelah semua rule dibuat, setiap cornerstone ditelusuri ulang. Semuanya tetap berakhir di
node yang dibuat untuknya, sehingga penambahan rule tidak merusak pengetahuan lama.

\begin{table}[H]
\centering
\begin{tabular}{llll}
\toprule
\textbf{Cornerstone} & \textbf{Jalur pada tree akhir} & \textbf{Output} & \textbf{Sesuai} \\
\midrule
C2 (R1) & R0(T), R1(T), R3(F), R4(F) & Saham & Ya \\
C3 (R2) & R0(T), R1(F), R2(T), R5(F) & Reksadana & Ya \\
C4 (R3) & R0(T), R1(T), R3(T) & Deposito & Ya \\
C6 (R4) & R0(T), R1(T), R3(F), R4(T) & Reksadana & Ya \\
C8 (R5) & R0(T), R1(F), R2(T), R5(T) & Obligasi & Ya \\
\bottomrule
\end{tabular}
\caption{Uji regresi cornerstone}
\end{table}

\subsection{Penanganan Conflict}
Kasus C9 identik dengan cornerstone R1 (C2), tetapi pakar meminta kesimpulan \emph{Obligasi}.
Tidak ada kondisi yang dapat membedakan keduanya, sehingga difference list kosong. Sistem
menampilkan peringatan dan tidak mengubah rule lama. Pakar dapat membatalkan koreksi, atau
menambah atribut baru, misalnya \texttt{tanggungan} (numerik) dengan nilai 2 untuk C9. Karena
C2 tidak memiliki nilai \texttt{tanggungan}, kandidat \cond{tanggungan >= 2} dan
\cond{tanggungan <= 2} bernilai FALSE untuk C2 dan dapat dipakai sebagai kondisi rule baru.
Alur ini diuji otomatis pada \texttt{tests/test\_cli.py}.

\subsection{Edge Case}
\begin{longtable}{P{0.38\textwidth}P{0.56\textwidth}}
\toprule
\textbf{Situasi} & \textbf{Perilaku sistem} \\
\midrule
\endhead
Input numerik bukan angka, di luar rentang, atau kategori di luar daftar & Ditolak dengan pesan
  kesalahan; pada CLI ditanyakan ulang. \\
Atribut belum diisi & Ditolak; semua atribut wajib diisi. \\
Pakar memilih kondisi yang juga dipenuhi cornerstone & Ditolak karena cornerstone ikut
  teranulir. \\
Kondisi kustom tidak dipenuhi kasus baru & Ditolak. \\
Kesimpulan koreksi sama dengan output sistem & Dianggap setuju, tidak ada rule baru. \\
Kesimpulan diketik dengan ejaan berbeda (``saham '') & Dinormalisasi menjadi kesimpulan yang
  sudah ada (``Saham''). \\
Difference list kosong & Conflict: batal atau tambah atribut (lihat subbab sebelumnya). \\
Urutan rule (C7) & R3 bernilai TRUE sehingga R4 di false-child R3 tidak dievaluasi. Ini
  perilaku yang benar pada SCRDR. \\
Kesimpulan default diubah & Aman karena root tidak memiliki cornerstone. \\
\bottomrule
\caption{Penanganan edge case}
\end{longtable}

\subsection{Pengujian Otomatis}
Proyek dilengkapi 52 unit test (pytest) yang seluruhnya lulus. Cakupannya:
\begin{itemize}[nosep]
  \item evaluasi kondisi numerik dan kategorikal, termasuk atribut yang tidak ada;
  \item pembangkitan difference list (dengan dan tanpa cornerstone, nilai sama, atribut baru);
  \item inferensi dan penempatan node pada true-child maupun false-child;
  \item validasi penambahan rule dan normalisasi kesimpulan;
  \item simpan, muat, dan reset KB serta validasi schema domain;
  \item skenario C1--C9 lengkap dengan bentuk tree akhir dan uji regresi;
  \item smoke test antarmuka Streamlit dan alur CLI (termasuk conflict dan rekaman log).
\end{itemize}

% ==================================================================
\section{Kesimpulan dan Saran}

\subsection{Kesimpulan}
Sistem pakar berbasis SCRDR berhasil dibangun sesuai spesifikasi. KB dimulai dari satu rule
default dan berkembang melalui koreksi pakar. Setiap rule baru adalah pengecualian bagi
kesimpulan yang salah, dengan kondisi yang dipilih dari difference list terhadap cornerstone.
Engine bersifat generik terhadap domain dan mendukung atribut numerik maupun kategorikal.
Pengujian skenario menunjukkan tidak ada regresi pada cornerstone lama.

\subsection{Saran}
\begin{itemize}[nosep]
  \item Mendukung \emph{Multiple Classification RDR} agar satu kasus dapat memiliki lebih dari
        satu kesimpulan.
  \item Menyimpan kasus yang disetujui di root sebagai cornerstone tambahan, agar rule
        pengecualian pertama tidak mengubah hasil kasus-kasus tersebut.
  \item Menyediakan nilai atribut baru untuk cornerstone lama saat atribut ditambahkan.
\end{itemize}

% ==================================================================
\section{Pembagian Kerja}
% TODO: sesuaikan.
\begin{table}[H]
\centering
\begin{tabular}{ll}
\toprule
\textbf{Anggota} & \textbf{Tugas} \\
\midrule
\textless Nama Anggota 1\textgreater{} (13523065) & \ldots \\
\textless Nama Anggota 2\textgreater{} (13523117) & \ldots \\
\bottomrule
\end{tabular}
\end{table}

\section*{Lampiran}
\begin{itemize}[nosep]
  \item Repositori: \textless tautan repositori\textgreater
  \item Contoh rekaman sesi CLI: \texttt{docs/contoh\_sesi\_cli.txt}
\end{itemize}

\end{document}
